\section{Tor网络攻防技术}

Tor（The Onion Router）是当前部署规模最大、研究最充分的低时延匿名通信系统之一。自第二代洋葱路由系统被提出以来，Tor逐步形成了由目录权威、志愿者中继、客户端、出口策略、洋葱服务与可插拔传输共同构成的复杂匿名通信生态\cite{dingledine2004tor,torspec}。从安全研究视角看，Tor并非单一密码协议，而是一个同时受到密码学设计、互联网路由、流量统计学习、系统资源管理与审查对抗共同影响的大规模网络系统。因此，围绕Tor的攻防研究也呈现出明显的多层化特征：攻击侧从早期的节点控制与流量关联，发展到深度学习网站指纹、AS/BGP级观测、洋葱服务会话关联、带宽膨胀和目录共识破坏；防御侧则从电路层协议加固、守卫节点策略和目录共识机制，扩展到流量填充、路由感知路径选择、隐私保护测量、抗审查传输和面向洋葱服务的性能优化。本节按照``架构基础--攻击演进--防御对抗--态势分析''的逻辑，对Tor网络攻防技术进行系统化梳理，并为后续讨论匿名网络与僵尸网络隐匿通信之间的关系提供技术基础。

\subsection{Tor网络架构}

\subsubsection{系统组成与节点角色}

Tor的思想源于一代洋葱路由。Syverson、Goldschlag与Reed提出的匿名连接模型已经奠定了``发送方构造多层加密洋葱、每个中继只剥离一层并获得下一跳''的基本范式\cite{syverson1997anonymous}。Dingledine等人提出的Tor则在此基础上完成了面向互联网部署的系统化重构：其核心目标不是提供绝对意义上的不可追踪性，而是在可接受时延和带宽开销下，使任一单点观察者难以同时获得通信双方身份\cite{dingledine2004tor}。这种设计决定了Tor本质上属于低时延匿名通信系统，适合Web访问、即时通信和一般TCP应用，但并不承诺抵御能够同时观察全网入口与出口流量的全局被动攻击者。

Tor网络由三类基本实体组成。第一类是\textbf{客户端}，负责获取目录共识、选择路径、建立电路并发起应用流量。第二类是\textbf{中继节点}，根据其在路径中的位置可分为守卫节点（Guard）、中间节点（Middle Relay）和出口节点（Exit）。守卫节点直接连接客户端，能够看到用户IP地址但不知道最终访问目标；出口节点负责将流量发送至公网目标，能够看到目的站点及未加密应用层内容但不知道用户真实地址；中间节点只知道前后两跳，用于切断入口与出口之间的直接关系。第三类是\textbf{目录基础设施}，包括目录权威与目录缓存。目录权威周期性投票生成网络状态共识，发布中继身份、公钥摘要、地址、端口、带宽权重及\texttt{Guard}、\texttt{Exit}、\texttt{HSDir}、\texttt{BadExit}等标志位，是客户端进行路径选择和抵御大规模Sybil渗透的基础\cite{torspec}。

\subsubsection{核心协议：目录共识、电路构建与洋葱加密}

Tor客户端在通信前首先下载并验证目录共识。共识文件给出了当前可用中继及其权重，客户端据此按照带宽加权、角色约束、出口策略和守卫节点策略选择一条典型的三跳路径。随后客户端采用增量式电路构建机制建立匿名通道：先与第一跳守卫节点完成认证密钥交换，再通过已经建立的前缀电路向第二跳、第三跳逐步发送扩展请求，使每个中继只获知自身相邻节点，而无法恢复完整路径\cite{dingledine2004tor,torspec}。这一\emph{telescoping construction}机制同时保证了路径隐私与每跳密钥独立性。

Tor的密码学结构可以概括为``非对称握手建立每跳密钥，对称加密承担高频数据转发''。在电路创建阶段，客户端分别与电路上的每个中继协商独立的会话密钥；在数据传输阶段，应用数据被切分为固定大小的信元并进行多层对称加密。客户端发送方向上，数据依次使用出口、中间、守卫三跳密钥进行嵌套加密；沿路径转发时，每个中继只解开属于自己的一层并根据内部头部信息转发给下一跳。反向传输时，出口、中间、守卫依次增加加密层，客户端最终逐层解开。固定大小信元、多路复用TCP流和端到端完整性校验共同降低了简单包长指纹和标签攻击风险，但并不能消除数据包方向、突发结构、时序间隔与连接持续时间等元数据泄露。

\subsubsection{洋葱服务与真实网络测量}

除访问公网服务外，Tor还支持\textbf{洋葱服务}（Onion Service）。洋葱服务允许服务端隐藏自身IP地址，并通过引入点与汇合点完成客户端和服务端的匿名连接。与普通出口电路相比，洋葱服务通常需要客户端侧三跳电路与服务端侧三跳电路共同参与，路径更长、交互轮次更多，因此在延迟、吞吐和可用性上面临更大压力。Winter等人通过访谈、问卷和\texttt{.onion}查询数据研究发现，用户对洋葱服务的心理模型并不完整，对地址发现和站点认证机制存在明显不满，同时现实中还存在大量相似洋葱地址和钓鱼风险\cite{winter2018tor}。这说明洋葱服务的安全问题不仅是协议问题，也包含可用性、身份认证和用户认知问题。

对真实Tor网络的理解依赖隐私保护测量。Mani等人基于PrivCount和Private Set-Union Cardinality对Tor进行测量，发现Tor日均用户规模可能显著高于传统估计，并观察到大量失败的洋葱地址查询\cite{mani2018understanding}。该类工作的重要意义在于：Tor研究不能只停留在协议规范或小规模实验网络中，而必须在不伤害用户隐私的前提下理解真实流量、真实服务和真实攻击面。后续网站指纹、AS级对手、洋葱服务性能和目录基础设施安全研究，均以这种测量基础为重要前提。

\subsection{攻击技术分类与演进}

针对Tor的攻击可按照攻击者观测面和攻击目标划分为四类：第一类是位于客户端侧或局部链路上的流量分析攻击，典型代表是网站指纹；第二类是同时观测入口侧与出口侧的去匿名化攻击，包括端到端流关联、DNS增强关联和AS/BGP级攻击；第三类是面向洋葱服务的服务发现、可用性和会话关联攻击；第四类是针对Tor生态基础设施的攻击，包括拒绝服务、带宽测量操纵和目录共识破坏。

\subsubsection{流量分析攻击：Website Fingerprinting的深度学习化}

网站指纹攻击（Website Fingerprinting, WF）试图在不破解Tor加密内容的情况下，仅利用包方向、突发长度、时间间隔和连接持续时间等元数据识别用户访问的网站或页面。2018年前后，WF研究发生了从手工特征工程向深度学习的转变。Sirinam等人提出Deep Fingerprinting（DF），使用卷积神经网络直接学习Tor访问轨迹中的方向序列特征，在未防御流量和部分轻量级填充防御下取得了远高于传统方法的准确率\cite{sirinam2018deep}。这一结果表明，固定大小信元只能削弱单包长度泄露，无法消除流量突发结构中蕴含的站点特征。

随后研究进一步关注模型的数据效率和特征完整性。Bhat等人提出Var-CNN，引入膨胀因果卷积和残差结构以捕捉长程时序依赖，并将自动学习特征与统计特征融合，证明深度WF攻击在低数据场景下仍具有较强能力\cite{bhat2018var}。Rahman等人的Tik-Tok则系统证明了时序信息并非噪声，而是与方向特征互补的重要信号；其方向--时序融合表示在未防御Tor、WTF-PAD和洋葱站点流量上均显著提升识别能力\cite{rahman2019tik}。到2024年，Zhao等人提出Oscar，将问题从`识别网站''推进到`识别同一网站内部具体网页''，通过多标签度量学习处理多页面混合和细粒度网页差异，在大规模监控页与非监控页数据集上显著提升Recall@5\cite{zhao2024towards}。这表明Tor流量分析的前沿已经从粗粒度站点分类转向开放世界、多标签、细粒度和跨时间泛化。

\subsubsection{去匿名化攻击：流关联、DNS泄露与AS级对手}

当攻击者能够同时观察客户端到守卫节点的入口侧流量和出口节点到目的站点的出口侧流量时，攻击目标便从`用户访问了哪类网站''升级为`某个用户是否正在访问某个具体目标''。Nasr等人提出DeepCorr，利用深度学习模型学习Tor流量在多跳转发中的噪声和扰动模式，从而对入口流和出口流进行端到端相关性判断\cite{nasr2018deepcorr}。该工作的重要性在于，它降低了流关联攻击对精心设计统计指标的依赖，说明中等规模多点观察者也可能借助数据驱动模型提升去匿名化能力。

DNS路径进一步扩大了流关联攻击的观测面。Greschbach等人的DefecTor研究指出，Tor出口节点对公共DNS解析器的选择存在集中化现象，攻击者若能观察出口DNS请求，便可将DNS信息与客户端侧网站指纹结合，从而显著提升关联攻击效果\cite{greschbach2016effect}。该工作提醒我们，Tor的匿名性不仅取决于Tor覆盖网络本身，还受出口侧DNS解析、目的站点部署和互联网基础设施集中化影响。

更底层的威胁来自AS级和BGP级对手。AS级攻击者不需要控制Tor中继，只需位于客户端--守卫路径和出口--目的路径的共同自治系统上，便可能被动关联两端流量。Sun等人提出Counter-RAPTOR，系统分析主动BGP前缀劫持对Tor匿名性的影响，说明攻击者可通过路由操纵使自身出现在入口或出口路径上\cite{sun2017counter}。Gegenhuber等人进一步利用RIPE Atlas长期测量不同国家客户端、Tor入口路径和目标AS之间的关系，展示了AS级流量关联风险的地域差异和带宽集中化效应\cite{gegenhuber2023extended}。因此，Tor威胁模型已经从单纯的覆盖网络内部节点控制，扩展到互联网路由与自治系统层面的联合分析。

\subsubsection{针对洋葱服务的攻击}

洋葱服务因同时隐藏客户端和服务端位置，长期被视为Tor中更强的匿名通信形态，但其更长路径和更复杂握手也引入了新的攻击面。Winter等人的用户研究表明，洋葱服务在地址发现、服务认证和用户信任方面存在天然脆弱性，攻击者可利用相似地址、过期链接和用户对\texttt{.onion}机制理解不足实施钓鱼或误导\cite{winter2018tor}。这类问题并不直接破坏Tor密码协议，却会在实际使用中削弱匿名服务的安全性。

在网络层，Lopes等人针对洋葱服务会话提出基于Sliding Subset Sum的流关联攻击，面向客户端与洋葱服务之间的会话流量进行过滤、时间对齐和滑动窗口匹配\cite{lopes2024flow}。与普通出口电路相比，洋葱服务路径更长、噪声更多、可能存在电路填充和多路复用，因此过去常被认为更难关联；该工作表明，只要攻击者在合适位置获得足够观测能力，仍可能通过高效统计匹配恢复会话关联关系。由此可见，洋葱服务安全不能仅依赖``路径更长''这一直觉，而需要结合流量填充、拥塞控制、服务发现和DoS防御进行系统设计。

\subsubsection{基础设施攻击与最新研究进展}

基础设施攻击关注的不是某条电路的瞬时匿名性，而是Tor网络维持可用性和可信路径选择所依赖的资源管理、带宽测量与目录共识。Jansen等人提出的Sniper Attack揭示了Tor早期拥塞控制和传输队列设计可被用于低成本拒绝服务：恶意客户端或服务端通过操纵读取行为和SENDME机制，使目标中继积累大量未释放数据，最终耗尽内存并崩溃\cite{jansen2014sniper}。该攻击的意义在于，DoS不只是可用性威胁，还可能通过选择性排除守卫或出口节点改变用户路径选择，进而服务于去匿名化。相关实验大量依赖Shadow仿真框架，Jansen等人关于Tor网络建模的方法论也成为后续安全评估的重要工具\cite{jansen2012methodically}。

2023年以来，研究焦点进一步转向带宽测量和目录共识机制。Sendner等人提出TorMult带宽膨胀攻击，利用多个中继共享资源以及带宽测量流量可被识别的缺陷，使攻击者控制的节点在共识中获得高于真实能力的带宽权重，从而增加被选入用户电路的概率\cite{sendner2023tormult}。这类攻击直接挑战了Tor``按带宽加权选择路径''的基础假设。进一步地，Luo等人针对目录协议提出短时DDoS破坏共识生成的研究，指出现有目录权威协议在网络同步假设上存在脆弱性，并提出基于部分同步模型的缓解方案\cite{luo2026five}。尽管该工作时间上已超出2023--2025窗口，但其揭示的方向代表了最新趋势：Tor攻防正在从客户端流量层面扩展到目录、测量和共识协议本身。

\subsection{防御技术与对抗}\label{sec:tor-defense}

Tor防御技术必须同时满足匿名性、性能、部署成本和兼容性约束。与高时延Mix网络不同，Tor不能通过长时间批处理、强制重排序或大量固定速率填充来彻底隐藏流量时序，否则会破坏Web浏览等交互式应用体验。因此，Tor防御研究的核心矛盾是：如何在有限带宽和低时延约束下，最大程度降低攻击者从元数据中提取稳定特征的能力。

\subsubsection{流量混淆技术}

针对网站指纹攻击，防御思路主要包括填充、延迟、突发整形和对抗扰动。Holland与Hopper提出RegulaTor，以下载流量速率衰减模型和上传/下载比例约束重塑网页加载轨迹，并在预算范围内注入哑包，从而降低WF分类器对突发形状的利用能力\cite{holland2020regulator}。其贡献不只在于降低分类准确率，更在于强调防御必须同时报告带宽开销、延迟开销和开放世界效果。Holland等人的DeTorrent进一步代表了填充型防御研究向对抗式、可部署方向发展的趋势，即在尽量不改变协议主体的条件下，通过更精细的padding策略降低流量分析收益\cite{holland2023detorrent}。

但是，流量混淆存在天然上限。若防御过弱，深度模型仍可利用剩余方向和时序模式；若防御过强，则会造成不可接受的带宽与延迟膨胀。更重要的是，攻击者可以采用自适应训练，将防御后的流量纳入训练集，使静态规则防御迅速失效。因此，当前更有前景的方向是将填充策略与真实网络状态、应用层加载阶段、拥塞控制和客户端部署环境结合起来，形成动态、低开销、可更新的防御机制，而不是依赖固定模式的离线混淆。

\subsubsection{协议层防御}

协议层防御主要包括守卫节点策略、路由感知路径选择、目录共识加固和隐私保护测量。守卫节点机制通过长期固定少量入口节点，降低用户在频繁随机选择入口时遇到恶意守卫的概率；目录权威和共识签名则通过多方投票降低单点目录投毒风险\cite{torspec}。针对AS/BGP级攻击，Counter-RAPTOR提出将AS弹性纳入守卫选择，并结合BGP异常监测对路由劫持进行检测\cite{sun2017counter}。这一思路说明，Tor协议层防御不能只优化覆盖网络内部路径，还必须考虑底层互联网路径是否会把入口侧和出口侧同时暴露给同一观察者。

面对基础设施攻击，防御重点则转向资源约束和共识鲁棒性。Sniper Attack之后，Tor需要在拥塞控制、传输队列和流控窗口上限制攻击者可放大的资源消耗\cite{jansen2014sniper}；TorMult之后，带宽权威测量需要更强的抗操纵能力，避免共享资源、测量流量识别和策略性丢弃用户流量造成路径选择偏斜\cite{sendner2023tormult}。目录协议方面，部分同步模型和更稳健的投票恢复机制可增强目录权威在短时DDoS或网络分区下的可用性\cite{luo2026five}。这些研究表明，Tor匿名性不仅依赖密码学正确性，还依赖测量系统、共识协议和网络资源治理的可信性。

\subsubsection{新兴防御技术}

新兴防御主要集中在两条路径：一是提升洋葱服务和Tor传输性能，二是增强抗审查能力。在性能方面，DarkHorse提出基于UDP的数据传输框架，试图缓解洋葱服务中TCP-over-TCP、长路径和多轮握手带来的高延迟问题\cite{al2023darkhorse}。该类方案的挑战在于，UDP带来的乱序、丢包和路径语义变化需要与Tor现有电路抽象、可靠性机制和匿名性假设重新协调。

在抗审查方面，Tor长期依赖可插拔传输将Tor流量伪装成更难被识别或更高封锁代价的协议流量。MacMillan等人对Snowflake的研究表明，基于WebRTC的抗审查传输虽然能利用大量临时代理提高可达性，但其握手特征和流量行为仍可能与真实WebRTC应用存在可区分差异\cite{macmillan2020evaluating}。Sun与Shmatikov提出的差异降级攻击进一步说明，审查者不一定需要准确分类流量，只需制造某些网络退化条件，使规避系统性能下降而正常应用仍可接受，即可形成低成本封锁\cite{sun2024differential}。Vilalonga等人的Obscura则代表了后续方向：将流量封装进WebRTC媒体流，结合瞬时代理和可靠传输层，试图同时抵抗深度包检测、代理枚举和差异降级\cite{vilalonga2026obscura}。总体看，抗审查研究已经从`让Tor不像Tor''转向`让封锁Tor的代价接近封锁正常应用''。

\subsection{攻防态势分析}

综合现有文献，Tor攻防态势呈现出四个趋势。第一，\textbf{攻击能力从规则化走向数据驱动}。DF、Var-CNN、Tik-Tok和Oscar表明，深度学习已经能够从加密流量的方向、时序和突发结构中自动提取稳定特征\cite{sirinam2018deep,bhat2018var,rahman2019tik,zhao2024towards}；Shadbeh等人的真实开放世界评估进一步说明，网站指纹并非纯粹实验室威胁，而可能在接近真实守卫观察者的条件下保持较高有效性\cite{shadbeh2026reality}。这迫使防御研究必须从封闭世界准确率转向开放世界精度、召回率、基准率、跨时间泛化和真实部署开销。

第二，\textbf{去匿名化攻击面从Tor覆盖网络内部扩展到底层互联网}。DeepCorr证明入口--出口流关联可由深度模型增强\cite{nasr2018deepcorr}，DefecTor说明DNS解析基础设施会扩大出口侧观测面\cite{greschbach2016effect}，Counter-RAPTOR和后续AS级测量则表明BGP路由、自治系统路径和中继带宽集中度都会改变用户实际暴露风险\cite{sun2017counter,gegenhuber2023extended}。因此，Tor匿名性不能仅由三跳电路结构解释，而必须放入互联网路由、公共解析器和跨域流量监测的环境中评估。

第三，\textbf{基础设施安全成为Tor生态的关键瓶颈}。Sniper Attack、TorMult和目录DDoS研究分别指向资源管理、带宽测量和目录共识三个核心支柱\cite{jansen2014sniper,sendner2023tormult,luo2026five}。攻击者若能影响这些基础设施，即使不破解任何密码原语，也可能通过改变路径选择概率、降低网络可用性或迫使用户连接恶意中继来削弱匿名性。这一趋势对于僵尸网络研究尤其重要：僵尸网络若利用Tor作为C2隐匿层，其稳定性同样依赖Tor中继生态和目录系统；反过来，针对Tor基础设施的攻击也可能影响恶意流量隐匿与防御监测的平衡。

第四，\textbf{防御研究越来越强调可部署性而非单一安全指标}。RegulaTor和DeTorrent体现了低开销抗WF防御的工程约束\cite{holland2020regulator,holland2023detorrent}，Counter-RAPTOR体现了路由安全与匿名熵之间的权衡\cite{sun2017counter}，DarkHorse和Obscura则分别从性能和抗审查角度探索Tor系统边界\cite{al2023darkhorse,vilalonga2026obscura}。总体而言，Tor未来演进很可能不会依赖某个单一突破，而是由路径选择、拥塞控制、流量填充、目录共识、抗审查传输和隐私保护测量的协同优化共同推动。

由此可见，Tor网络的核心矛盾始终是低时延可用性与强匿名性之间的张力。攻击者只需在入口观测、出口观测、路由操控、流量统计或基础设施治理中的某一环节取得优势，即可显著削弱匿名性；防御者则必须在不牺牲可用性和开放部署生态的前提下，同时约束多类攻击面。这种不对称性决定了Tor仍将是匿名通信、安全测量和攻防博弈研究中的长期核心对象，也为后续分析僵尸网络如何借助匿名网络隐藏指挥控制、以及安全防御方如何在不侵犯正常用户隐私的条件下识别滥用行为，提供了直接的理论与技术铺垫。
